% Generated draft: requires mathematical and source review.
\documentclass[11pt]{article}
\usepackage[margin=25mm]{geometry}
\usepackage{fontspec}
\setmainfont{Latin Modern Roman}
\usepackage{amsmath,amssymb,mathtools,unicode-math}
\setmathfont{Latin Modern Math}
\usepackage{enumitem,needspace,xurl,hyperref}
\setlength{\parindent}{0pt}
\setlength{\parskip}{5pt}
\setlength{\emergencystretch}{3em}
\newcommand{\N}{\mathbb{N}}
\newcommand{\Z}{\mathbb{Z}}
\newcommand{\Q}{\mathbb{Q}}
\newcommand{\R}{\mathbb{R}}
\newcommand{\C}{\mathbb{C}}
\newcommand{\F}{\mathbb{F}}
\newcommand{\A}{\mathbb{A}}
\newcommand{\PP}{\mathbb{P}}
\newcommand{\E}{\mathbb{E}}
\newcommand{\eps}{\varepsilon}
\newcommand{\dd}{\,\mathrm{d}}
\begin{document}
\section*{UnsolvedMath ID: 30003509}
\textbf{Source problem code:} catalogue entry
\subsection{Definitions and mathematical statement}
A number field is a finite field extension \(F/\Q\). For each number field \(F\), let \(\mathcal R_F\) denote the class of reductive algebraic groups over \(F\) intended by the cited source. Write \(\operatorname{TWN}^{+}(G/F)\) for the assertion that \(G\) satisfies the source's property \((\mathrm{TWN}+ )\); this notation names that property and does not define it.

\textbf{Conjecture (as supplied).} For every number field \(F\) and every \(G\in\mathcal R_F\), the group \(G\) satisfies property \((\mathrm{TWN}+ )\). Symbolically,
\[
\forall F\text{ a number field},\qquad
\forall G\in\mathcal R_F,\qquad
\operatorname{TWN}^{+}(G/F).
\]

\textbf{Unresolved formulation qualifications.} The supplied record does not define property \((\mathrm{TWN}+ )\), its analytic objects, or its required estimates and quantifiers. It also does not specify the connectedness convention for reductive groups. Consequently, the display preserves the source's universally quantified conjecture but is not a self-contained formulation of its analytic content. The precise definition and group conventions require checking in the supplied references; no replacement definition is asserted here.
\subsection{Short English statement}
The conjecture asks whether every reductive group over every number field has property \((\mathrm{TWN}+ )\). The supplied record omits the definition of that property, so its precise analytic requirements remain unresolved. The record labels the problem partially solved; that label and the accompanying literature assessment are supplied curation, not independently verified conclusions.
\subsection{Sources}
\begin{itemize}
\item Supplied local UnsolvedMath registry record: permanent ID \(30003509\), title \emph{Property (TWN+) for Reductive Groups}, exact source problem code \texttt{catalogue entry}, in the Oberwolfach Reports Open Problems collection. The record gives proposed year \(2017\) and names no individual proposer.
\item Supplied source citation: \emph{Harmonic Analysis and the Trace Formula (2018)}, Oberwolfach Reports, pp.~1551--1630, DOI \texttt{10.4171/owr/2017/25}. \url{https://doi.org/10.4171/owr/2017/25}. The background describes these reports as documenting mathematical workshops and their open-problem discussions.
\item Supplied literature reference: Finis and Lapid, \emph{analytic properties of intertwining operators}. \url{https://arxiv.org/abs/1603.05475}.
\item Supplied literature reference labelled \emph{original OWR report}. \url{https://doi.org/10.4171/owr/2017/25}.
\item Supplied literature reference labelled \emph{current Eisenstein-series bounds}. \url{https://doi.org/10.1017/fms.2024.59}.
\item Literature assessment dated exactly \texttt{2026-08-21}: the supplied curation labels the problem partially solved and states that the conjecture for all reductive groups remains unproved, while property \((\mathrm{TWN}+ )\) is established for \(\mathrm{GL}_n\) and inner forms, quasi-split classical and similitude groups, \(G_2\), and groups related by derived isogeny. This is supplied dated curation derived from the OWR record, not an independently verified assessment; its time-sensitive status and cited claims require independent verification.
\item The registry assigns the statement status \texttt{corrected\_verified} and reports that the official Oberwolfach report was checked directly, with definitions restored and source-dependent wording removed. These are supplied verification claims. No source visit or independent verification is claimed in this draft.
\end{itemize}
\end{document}
